% Chapter 1 -- translation by Doug Burkett & Tim Hutcheson (ca. 2001),
% lightly copy-edited for this edition.

\chapter{General Symbulator Concepts}

\section{What is Symbulator?}

Symbulator is a linear circuit simulator which finds symbolic and
numeric results. Symbulator can perform direct current (DC) and
alternating current (AC) analyses. Symbulator can also find
frequency domain response and transient response. It supports a wide
variety of linear circuit elements. It occupies less than 60KB of
calculator memory and is programmed in BASIC. Symbulator runs on
particular Texas Instruments graphing calculators.

\section{Requirements}

Symbulator runs on the Texas Instruments TI-89 or equivalent
calculators which have AMS version 2.0x or superior installed. AMS
is an acronym for Advanced Mathematics System; this is what Texas
Instruments calls the calculator operating system. You must also
have the program DiffEq (version 2.06 or superior) installed. Both
Symbulator and DiffEq must be archived. In the rest of this book, I
refer to a calculator that meets these requirements as \emph{your
calculator}.

\subsection{Installation information resources}

If your calculator does not meet these requirements, you must update
it before you can use Symbulator. As this book focuses on using
Symbulator, and not other aspects of calculator operation, I will
not provide upgrade details. The following resources will help you
upgrade your calculator.

To upgrade your calculator to AMS 2.05, refer to this Texas
Instruments Internet site: \code{http://education.ti.com}.

The Symbulator and DiffEq programs can be downloaded free from the
official Symbulator site shown in section 1.3 below. In order to
install these programs on your calculator, you will need both the
GraphLink software and a GraphLink cable to connect your calculator
to a personal computer. The GraphLink software is free and can also
be downloaded from \code{http://education.ti.com}.

\subsection{Archive information resources}

Symbulator and DiffEq execute much more quickly when they are
correctly archived in your calculator's flash memory. Both
Symbulator and DiffEq are distributed with programs which
automatically perform the archive operation. Instructions to use the
archiving programs are provided in the documentation for Symbulator
and DiffEq, which is included in the downloaded files.

\section{Symbulator Internet site}

Symbulator has an official web site at
\code{http://paxm.org/symbulator}. This site is in English. In
addition to complete documentation and many step-by-step example
problems, the most recent versions of Symbulator and DiffEq can be
downloaded free.

\section{Brief history of Symbulator}

I am Roberto P\'erez-Franco, the author of Symbulator. I am a
student of Electromechanical Engineering in the Technological
University of Panama, located in Azuero. I started programming
Symbulator, then called SCS, on March 30, 1999. The first released
version was 0.92, which was posted on the Symbulator web site on
April 2, 1999. Version 1 added mutual inductances. Version 2 was
released August 19, 1999, at Azuero and added support for two-port
circuits. Version 3 incorporated a new Expert system. Version 4 was
released on the Symbulator web site on January 9, 2001, and added
the Impala enhancement. The most recent version, called
Symbulator~Q, was released on the Symbulator web site on March 30,
2001. Symbulator~Q includes significant improvements which required
modifying about 40\% of the code.

Symbulator has always been offered for free on the Internet.
Feedback from users has allowed me to polish the code, eliminate
errors and add novel features.

Now, two years after its public appearance, Symbulator is used by
thousands of electrical engineering students and professionals
around the world, and has received several honors, including 1st
place in Student Projects of the IEEE (Institute of Electrical and
Electronics Engineers), for the year 2000 in Latin America
(Region~9).

\subsection{Language}

Spanish is my native language. I have several books published in
Spanish, and I am considered a promising young Panamanian author. In
addition to Spanish, I also speak English and Esperanto. I have also
studied French and wish to continue developing my literary skills.

When the Symbulator project began in 1999, TI-89/92+ calculators
were practically nonexistent in Panama. Users play a vital role
during the development of a program, both in finding bugs and
inspiring the author with comments, encouragement and thanks. I have
received no monetary compensation for this program; it has always
been free.

Since most users of the TI-89 and TI-92+ speak English as a first or
second language, it was appropriate to use that language for the
program interface and documentation. This made it easy for users on
five continents to send comments and suggestions to me. As the
capabilities of Symbulator extended in many directions in the last
two years, the rapid release of new versions left me little time to
keep the documentation up to date. For this reason the documentation
and the program interface continued to be written in English.

I offer my apologies to victims of this situation. It has always
been my desire to develop a useful tool for all the engineering
students of the world. In fact, my interest in Esperanto has its
root in the lack of a true neutral, common international language
for all people.

Note that I have responded in Spanish to all users who used that
language in correspondence with me. Also, this book---which is the
most complete documentation of Symbulator---was originally written
in Spanish. You are reading the English translation, and I plan to
translate this book to Esperanto and other languages as soon as
possible.

\section{Concept of the algorithm}

In order to analyze an electrical circuit, Symbulator uses methods
that emulate the manual techniques of circuit analysis which are
learned by students. In fact, Symbulator operates like an excellent
student who has paid attention and learned his lessons well. As the
program uses the same methods as a student, it is also a very
powerful tool for learning circuit analysis.

Symbulator knows the Ohm--Cavendish law and the Kirchhoff current
law. With these two laws, the program analyzes circuits using only
nodal analysis. Although Kirchhoff's voltage law is not used, it is
always satisfied as a consequence of the other two laws. Mesh
analysis is not used. Thanks to the powerful built-in capabilities
of the calculator, nodal analysis is sufficient, even for the most
complex circuits.

Symbulator also knows the Th\'evenin--Helmholtz and Norton--Mayer
theorems. Symbulator can also graph Bode plots. To perform the
necessary Laplace transforms and inverses, Symbulator uses the
DiffEq program. There is no faster program for Laplace transforms
and inverses for any computer or calculator. DiffEq is the fruit of
the brilliant mind of Lars Frederiksen, whom the author is proud to
count among his friends.

Symbulator generates the equations of each node and each voltage
source, as well as the equations for each special element. At the
same time, Symbulator identifies the primary circuit unknowns.
Finally, it finds the symbolic or numeric value of all the primary
unknowns which solve the equations, and then finds the values of the
secondary and tertiary unknowns from the primary unknowns. Thus, it
offers all the circuit information a student could wish.

\section{Doors}

The heart of Symbulator consists of several programs called
\code{step} programs, which actually perform the simulation. The
user never directly uses these \code{step} programs, but instead
uses other programs that act as \emph{doors}, which in turn use the
\code{step} programs. There are four access doors, one for each
basic type of analysis: 1) \code{dc} is used for direct-current
analysis, 2) \code{ac} is used for alternating-current analysis, 3)
\code{fd} is used for frequency domain analysis, and 4) \code{tr} is
used for transient analysis.

\section{Tools}

Symbulator includes eight more tools that can help with performing
simulations or understanding the simulation results: 1) \code{par}
reduces parallel resistances to a single resistance, 2)
\code{absang} shows complex answers in polar form, 3) \code{plot}
graphs a function of time, 4) \code{bode} draws Bode plots, 5)
\code{thevenin} finds the Th\'evenin equivalent of a circuit, 6)
\code{port} finds the equivalent parameters for a two-port circuit,
7) \code{gain} finds the gain of an amplifier, and 8)
\code{fd\_to\_tr} converts a frequency domain response to the time
domain. Throughout this book you will see how these tools are used,
as well as the tools' syntax.

For quick reference, Symbulator includes a concise Help function. To
display the Help screens, enter \code{sq\char`\\help()} in the entry
line. The Help is in English.

\section{Commands}

Sometimes to solve a problem with Symbulator, you must use the
built-in calculator commands, such as those which solve equations.
You may also need to use auxiliary programs such as the Laplace
functions of DiffEq.

\section{Circuit descriptions}

\subsection{Philosophy}

The philosophy of the Symbulator circuit description is the
principle that you need only provide the minimum necessary
information to completely describe each circuit element.

\subsection{Text description}

Since Symbulator is designed to be used quickly in engineering
classes, circuits are described as text strings. A graphical
interface is not used, because that would tax the calculator's
processor, which is a Motorola 68000 running at 12~MHz. Also, this
method would make entering the circuits slow.

A circuit is described as a text string. Each \emph{row} of the
string describes one circuit element. Each row almost always has
five columns. The first element of each row is the element name. The
element name must begin with a letter that identifies the type of
circuit element the row describes. The next row elements are the
node names, which specify where the circuit element is connected.
The node names may be numbers or letters. The next row element is
the value of the circuit element in its corresponding unit of
measurement. The element value may be a number, a symbolic variable
name, or even an algebraic expression. For elements which store
energy (inductors and capacitors), their initial conditions are
placed in the last row location.

The elements in each row are separated by commas. The rows are
separated by semicolons. This example shows a circuit with three
circuit elements, and five terms for each element:

\begin{entry}
a1,b1,c1,d1,e1;a2,b2,c2,d2,e2;a3,b3,c3,d3,e3
\end{entry}

Internally, Symbulator converts this string to a matrix, to work
more quickly with the circuit description.

\subsection{Nodes}

\subsubsection{Order of the nodes}

The order in which the node names are specified in the circuit
description is not arbitrary: the order of the node names indicates
the direction of current flow and voltage drop in the circuit
element. This is particularly important for voltage and current
sources, for elements which control dependent sources, for elements
which store energy and have initial conditions, and for mutual
inductances.

\subsubsection{Names of the nodes}

Before you can run the simulation, you must identify each circuit
node and element with a unique name.

In the case of node names, the name can be a number, letter or a
combination of both. Each node has only one name, which means you
must use that same name every time the node is used in the circuit
description.

If you are not accustomed to using circuit simulators, you may
forget the node name once the simulation results are shown. It may
help, as you learn to use Symbulator, to write the node names and
component names on the circuit schematic.

\subsubsection{Reference node}

Each circuit always has exactly one reference node. The voltage at
the reference node is always 0 volts. This node is known as ground,
or earth, and must always be given the node name 0 (zero).

\subsection{Reserved variables}

The calculator operating system reserves some variable names for its
own use. These variables are called reserved or system variables.
The TI-89/TI-92+ User's Manual lists these reserved variable names.
You cannot use reserved variable names as component or node names.
In particular, be sure that you do not use these names: \code{nc},
\code{ok}, \code{rc}, \code{sx}, \code{tc}, \code{yc}, \code{zc},
and \code{c1} to \code{c99}.

Symbulator also uses many internal variables to perform the
simulation. There may be many of these variables, and they are
generated according to the names of the circuit elements, so it is
impossible to specify these names. However, they all begin with
Greek letters, so it is easy to avoid accidentally using these
names: do not use component names or node names that start with
Greek letters.
